\documentclass[11pt]{article}
\usepackage[a4paper,margin=28mm]{geometry}
\usepackage{amsmath,amssymb,amsthm}
\usepackage[hidelinks]{hyperref}
\newtheorem{theorem}{Theorem}
\title{The Satisfiability Hypothesis from fixed-power graph bounds}
\author{Atlas\thanks{The proof in this note was found using ChatGPT.}}
\date{18 November 2026}
\begin{document}
\maketitle

We consider finite simple undirected graphs with individually represented
vertices and a nonnegative integer threshold \(k\) in binary. Let \(n\)
denote the number of vertices and \(B\) the complete input length in bits.
In \textsc{Independent Set} the question is whether at least \(k\)
pairwise nonadjacent vertices exist. In \textsc{Vertex Cover} the question
is whether at most \(k\) vertices suffice to meet every edge.
The Satisfiability Hypothesis (SH) excludes deterministic CNF satisfiability
algorithms in time \(2^{O(s^a)}s^{O(1)}\) for each fixed \(0<a<1\), where
\(s\) is the normalized input length~\cite{sh}.

\hypertarget{is-order-sh-converse}{}
\hypertarget{vc-order-sh-converse}{}
\hypertarget{vc-budget-sh-converse}{}
\begin{theorem}\label{thm:classical-sh}
Each of the following exclusions implies SH. In each exclusion \(a\)
ranges over all fixed values \(0<a<1\).
\begin{enumerate}
\item No deterministic \(2^{O(n^a)}B^{O(1)}\)-time algorithm for
\textsc{Independent Set}.
\item No deterministic \(2^{O(n^a)}B^{O(1)}\)-time algorithm for
\textsc{Vertex Cover}.
\item No deterministic \(2^{O(k^a)}B^{O(1)}\)-time algorithm for
\textsc{Vertex Cover}.
\end{enumerate}
\end{theorem}

\begin{proof}
Suppose SH fails, and fix a deterministic CNF solver of time
\(2^{O(s^b)}s^{O(1)}\) for some \(0<b<1\).
We first obtain an algorithm whose exponent depends on the number of
variables in a width-three formula. By quantitative
sparsification~\cite[Lemma~2.2]{sh}, there is a constant \(D\ge1\) such
that, for each fixed \(0<d<1\), a width-three formula on \(v\) variables
can be expressed as a disjunction of at most \(2^{O(v^{1-d})}\)
width-three formulas, each with at most \(v\) variables and
\(O(v^{1+Dd})\) clauses. Generation takes deterministic time
\(2^{O(v^{1-d})}\) times a polynomial in the complete formula input
length. The width-three sparsification statement is a preprint dependency.
Relabel the variables of each output formula consecutively so that its
encoding uses \(O((v+1)^{1+Dd}\log(v+2))\) bits.
Choose \(d\) small enough that \(b(1+Dd)<1\) and fix \(a\) strictly between
\(\max\{1-d,b(1+Dd)\}\) and \(1\). The strict inequality absorbs the
logarithmic factor in the encoding size. Running the CNF solver on
every output formula, and taking the disjunction, gives a deterministic
width-three solver with time \(2^{O(v^a)}\) times a polynomial in its
complete input length.

For the first two exclusions, we encode the graph problem in width-three
CNF on \(O(n+1)\) variables. Thresholds larger than \(n\) and the case
\(n=0\) admit direct decisions. Use one selection bit \(x_u\) per vertex.
For \textsc{Independent Set}, add
\(\neg x_u\lor \neg x_w\) for every edge \(uw\).
For \textsc{Vertex Cover}, add \(x_u\lor x_w\) for every edge.
The respective constraints on the number of selected vertices are
enforced by a Boolean circuit summing the bits and comparing their sum
with \(k\). Pad the bits to a power of two and sum them with a balanced
tree of binary adders. Level \(j\) uses \(O(n/2^j)\) additions of
\(O(j)\)-bit numbers. With ripple-carry adders, the whole circuit has
\(O(n\sum_{j\ge1}j/2^j)=O(n)\) constant-fanin gates, including the
final comparison. For every fanin-two gate use one variable for its
output and clauses forbidding each input-output combination inconsistent
with the gate. These clauses have width at most three. Force the final
comparison output to be true. The formula is satisfiable exactly when
the graph instance is a yes-instance, and its length and construction
time are polynomial in \(B\). The width-three solver now runs in time
\(2^{O(n^a)}B^{O(1)}\), contradicting either graph-order exclusion.

For the third exclusion, we first use the polynomial-time
\(2k\)-vertex kernel for \textsc{Vertex Cover}, due to Nemhauser and
Trotter~\cite{nt,jb}. It either decides the input directly or produces an
equivalent graph with at most \(2k\) vertices and an updated threshold.
The same encoding uses \(O(k+1)\) variables and has length polynomial in
\(B\). Bounded \(k\) can be handled in polynomial time on the kernel.
The resulting running time is \(2^{O(k^a)}B^{O(1)}\), contradicting
the budget exclusion.
\end{proof}

\begin{thebibliography}{9}
\bibitem{sh}
MIT Hardness Group, J.~Brunner, E.~D. Demaine, F.~Frei, T.~Gomez,
M.~Hecher, J.~Li, and J.~Lynch.
\emph{ETH is Dead, Long Live SH}.
\href{https://arxiv.org/abs/2611.02650v1}{arXiv:2611.02650v1}, 2026.

\bibitem{nt}
G.~L. Nemhauser and L.~E. Trotter.
Vertex packings: Structural properties and algorithms.
\emph{Mathematical Programming}, 8:232--248, 1975.
\href{https://doi.org/10.1007/BF01580444}{DOI}.

\bibitem{jb}
B.~M. P. Jansen and H.~L. Bodlaender.
Vertex Cover Kernelization Revisited.
\emph{Theory of Computing Systems}, 53:263--299, 2013.
\href{https://doi.org/10.1007/s00224-012-9393-4}{DOI}.
\end{thebibliography}
\end{document}
